\section{Conclusions}
\label{sec:conclusions}

In this investigation, a two-dimensional Bloch--Floquet finite element framework was developed for Mindlin Form-II strain-gradient elasticity incorporating an anisotropic characteristic-length tensor and micro-inertia. The principal conclusions are summarized as follows:
\begin{enumerate}
  \item \textbf{Anisotropic Length-Scale Formulation:} By deriving the characteristic length tensor $\mathbf{L}(\theta)$ from the second moments of an ellipsoidal averaging domain, microstructural aspect ratio $\mathrm{AR}$ and orientation $\theta$ were successfully incorporated as continuous, independent design knobs. The positive definiteness of $\mathbf{L}(\theta)$ was shown to be strictly invariant under spatial rotation.
  \item \textbf{Consistent $C^1$ Bloch Discretisation:} A 32-degree-of-freedom Bogner--Fox--Schmit bicubic rectangular element was formulated. We demonstrated that conforming $C^1$ representations require identical complex Bloch phase transformations across displacement, first-order, and cross-derivative nodal degrees of freedom, preserving the Hermiticity and periodicity of the reduced eigenproblem.
  \item \textbf{Rigorous Numerical Verification:} The finite element framework satisfied an eight-test internal consistency suite with zero violations. Monotone mesh convergence across $4^2$ to $32^2$ grids confirmed an empirical least-squares rate of $p = 4.17$ and an operational numerical resolution floor of $\varepsilon_\Delta = 4.63 \times 10^{-11}$.
  \item \textbf{Directional Band-Gap Filtering:} In homogeneous media (Case H), microstructural anisotropy lifts high-symmetry mode degeneracies, producing directional stop bands up to $\Delta_{GX} = 0.0896$ along $\Gamma$--$X$ while strictly adhering to the absence of complete band gaps ($\Delta_{\mathrm{complete}} \le -0.3758$). The orientation sensitivity metric $S_\theta$ increases monotonically with aspect ratio up to $3.946\,\mathrm{rad}^{-1}$ at $\mathrm{AR} = 10$.
  \item \textbf{Complete Omnidirectional Band Gaps via Periodic Heterogeneity:} Embedding high-contrast circular inclusions (Case~C) opens a broad complete omnidirectional band gap ($\Delta_{\mathrm{complete}} = 2.5732$, $\Delta/\bar{\omega}_{\mathrm{mid}} = 43.94\%$), strictly obeying the directional-to-complete subset hierarchy $\Delta[\mathrm{leg}] \ge \Delta[\mathrm{path}] \ge \Delta[\mathrm{complete}]$. The gap remains open under mesh and quadrature refinements while its width decreases with mesh refinement ($\Delta_{\mathrm{complete}} = 2.5732$, $2.2504$, $2.0722$ at $4\times 4$, $8\times 8$, $16\times 16$; not mesh-converged), so $\Delta_{\mathrm{complete}} = 2.5732$ is stated at the $4\times 4$ mesh.
  \item \textbf{Acoustic Wave Steering:} At constant wavenumber $|\bmk| = 0.5$, group velocity vectors deviate from phase vectors by up to $\delta_{\max} = 2.79^\circ$ at $\mathrm{AR} = 10$, demonstrating active acoustic steering without material contrast.
  \item \textbf{Necessity of Micro-Inertia:} Both analytical asymptotic derivations and numerical computations confirmed that omitting micro-inertia ($\ell_{\mathrm{i}} = 0$) results in an unbounded phase velocity ($v_p \propto k \to \infty$), whereas $\ell_{\mathrm{i}} > 0$ bounds the wave speed to a finite asymptotic horizon ($v_{T,\infty} = 0.3162$), ensuring dynamic causality.
\end{enumerate}

These findings provide a solid theoretical and computational foundation for designing architected metamaterials that leverage microstructural orientation for directional wave guiding and vibration suppression.
